\chapter*{How to use this book}
\phantomsection
\addcontentsline{toc}{chapter}{How to use this book}
\markboth{HOW TO USE THIS BOOK}{How to use this book}

\noindent
This book cannot be read. It has to be worked. That is not a figure of speech
and it is not encouragement: the format is built so that passive reading
produces almost nothing, and four pages in you will feel it.

\section*{The frame}

Each program is a stream of small numbered chunks called \emph{frames}. Nearly
every frame ends by asking you for something \dash{} a number, a word, a line of
working. \textbf{The answer is at the top of the next frame.}

\begin{warning}
Cover the next frame before you read the current one to the end. A sheet of
paper is enough. Write your answer down \dash{} on paper, not in your head \dash{}
and only then uncover.

This is the one rule that is not negotiable. Break it and the book degenerates
into a mediocre collection of worked examples, and every advantage the method
has is gone. You will be tempted to break it in exactly the frames where it
matters most, which are the ones where you are not sure.
\end{warning}

\noindent
A row of dots \blank{} means a word, a number or an expression is missing and
you are to supply it.

\section*{The scaffolding}

Inside one program the examples move through four levels, in order, and the
change is deliberate:

\begin{enumerate}[leftmargin=1.6em, itemsep=3pt]
  \item an example worked in full, with the reasoning shown;
  \item an example worked with gaps for you to fill;
  \item \emph{now one for you to do} \dash{} you solve it, and check in the next
        frame;
  \item the end-of-program problems, which you do entirely unaided.
\end{enumerate}

\noindent
By the time you reach level four you have done the technique three times with
decreasing help, which is why it tends to feel easy rather than exposed. That
effect is the point of the design.

\section*{The seven parts of every program}

\begin{center}
\small
\begin{tabularx}{\linewidth}{@{}lX@{}}
\toprule
\textbf{Learning outcomes} & What you will be able to do, listed before you
  start, so you can judge the program against its own promise. \\
\textbf{Quiz} & Take it \emph{before} the program: it tells you whether you
  need the program at all, and if you only need part of it, which part. Each
  question names the frames that teach it. It is a diagnostic, never a score. \\
\textbf{Frames} & The program itself. The Summary and the Test exercises
  are numbered frames too, and they carry the last numbers in the program. \\
\textbf{Summary} & Every item tagged with the frame it came from,
  \mbox{like [17]}. The summary is therefore also a route back. \\
\textbf{Can you?} & The learning outcomes again, one for one, to rate yourself
  against from 1 to 5. \\
\textbf{Test exercises} & Exactly what the program taught, in the order it
  taught it. No traps. \\
\textbf{Further problems} & Consolidation, and more of it than you will want. \\
\bottomrule
\end{tabularx}
\end{center}

\section*{The loop that actually works}

\begin{center}
\begin{tcolorbox}[enhanced, sharp corners, width=0.94\linewidth,
  colback=mlight, colframe=mblue, boxrule=0pt, leftrule=3pt,
  left=10pt, right=10pt, top=8pt, bottom=8pt]
read the outcomes $\rightarrow$ take the \textbf{Quiz} $\rightarrow$ work the
program $\rightarrow$ \textbf{Can you?} $\rightarrow$ \textbf{Test exercises}
$\rightarrow$ Further problems
\end{tcolorbox}
\end{center}

\noindent
Most people skip the Test exercises, and the Test exercises are the most
informative thing in the sequence. The Quiz is a diagnostic on the way in and
nothing else: answered a second time it is the same questions put to a reader
who now remembers the answers, so the difference between two Quiz scores
measures your memory. The Test exercises score what the program taught, and
what they report is frequently below the feeling of having understood it.

\section*{If you already have a degree}

Do not read Part~I in order. It assumes nothing, deliberately, and if you can
already differentiate a product you will find it slow.

\textbf{Work the thirteen Foundation Quizzes and nothing else.} They take an
afternoon. Enter a Foundation program only where its Quiz goes badly, and enter
it at the frames the failed question names rather than at frame~1. Then start
at Program~1 and read the main parts in order \dash{} they are not independent, and
the dependency list in the introduction says which programs assume which.

Two Foundation programs are worth working even if the Quiz goes well:
Program~F3 on logarithms, because almost everything numerical in this book
happens in log-space and the reasons are not the ones taught at school; and
Program~F12 on the chain rule, because backpropagation is the chain rule and
nothing else, and the book leans on that sentence for four hundred pages.

\section*{What this book claims, and what it does not}

\begin{note}
Stroud's programs were validated to an \textbf{80/80} standard before
publication: at least 80 per cent of students scoring at least 80 per cent. The
finding worth having is not that the mean rose but that the \emph{spread
narrowed} \dash{} the method lifts the bottom of the distribution rather than
the top.

\textbf{This book has not been tested on a single reader.} It borrows the
method, and it is not entitled to borrow the evidence. Every build of it prints
that as outstanding work, and it will keep printing it until somebody runs the
trial.
\end{note}

\begin{warning}
\textbf{And the general evidence for programmed instruction is not good.}
Kulik, Cohen and Ebeling's 1982 meta-analysis of the format in secondary
education found that it did not typically raise achievement, and did not make
students feel more positively about the subject. Small steps and linear
sequencing did not turn out to be essential; immediate reinforcement did not
turn out to be critical.

That is a finding against the format as a whole, and it is quoted here rather
than left out.

The reason this book uses the format anyway is narrower and better supported
than \enquote{programmed learning works}. It is \emph{retrieval practice}: with
total study time held constant, being made to recall something produces
markedly better long-term retention than rereading it. And it is
\emph{errorful generation}: eliciting a wrong answer and then correcting it
beats studying the same material without the error \dash{} with the benefit
largest exactly when the learner was confident. Those two findings are what the
frame and the trap are for. They are not a defence of small steps for their own
sake, and this book does not offer one.
\end{warning}

\begin{warning}
\textbf{It will feel worse than reading while you are doing it.} That is not a
figure of speech either. In the retrieval-practice studies the advantage is
absent at an immediate test and appears after two days and a week; readers
consistently rate the method they retain \emph{less} from as the more
effective one, because fluency feels like learning and effort feels like
failure.

If you abandon this book because it feels inefficient, you will be abandoning
it for the reason it works.
\end{warning}
